\section{Introduction}
\label{sec:intro}

Minimum-edge-cut recursive bisection produces congressional district maps by
repeatedly splitting the state's census-tract adjacency graph into two
population-balanced halves, minimising the total shared boundary length at each
level. The algorithm makes no reference to partisan data: it sees only
geography and population. Yet on states with concentrated metropolitan areas,
it reliably produces a specific pattern:

\begin{quote}
\emph{The first bisection peels off the densest urban core as a single compact
district, leaving the rest of the state to be bisected at subsequent levels.}
\end{quote}

We call this \textbf{urban peeling}. It arises not from any partisan intent
but from a geometric fact: the boundary surrounding a dense city is shorter
than any boundary that bisects the state into two comparably-populated regions.
Illinois splits 1:17 at the first level (Chicago peel), not 8:9.
New York splits 1:25 (New York City peel), not 13:13.
Wisconsin splits 1:7 (Milwaukee peel), not 4:4.

GeoSection (B.8) adds isoperimetric normalisation to prevent the worst form
of this pathology --- the caterpillar, where 1:k-1 wins at every single
level regardless of state shape --- but it does not eliminate urban peeling
for states where the urban core is genuinely compact.
A normalised 1:7 split in Wisconsin still wins when Milwaukee is as compact
as it is: the normalisation says ``this boundary is as short as a 1:7 boundary
could be,'' which does not tell us whether it is shorter than a 4:4 boundary.

Urban peeling is legally contentious for two reasons.
First, it concentrates a disproportionate share of the minority party's voters
into a single safe district at every recursion level, producing implicitly
packed maps.
Second, it creates near-toss-up districts in the suburban ring around the
peeled city --- a predictable byproduct of the boundary geometry --- that
can influence which party controls the chamber.
Both effects arise with zero partisan input, but they are not politically
neutral.

\paragraph{Our contribution.}
We introduce \textbf{AreaSection}, which adds a second METIS balance constraint
requiring each half of every bisection to receive approximately 50\,\% of the
state's land area, with a $\pm 10$\,\% swing tolerance.
The area constraint is enforced simultaneously with the population constraint
using METIS's native dual-constraint (\texttt{ncon=2}) partitioning with
per-constraint imbalance vectors (\texttt{ubvec}).

The area constraint cannot be satisfied by a 1:17 split in Illinois because
Chicago occupies $<5$\,\% of the state's land.
A near-equal population split is therefore geometrically forced by the area
balance requirement.
Illinois moves from 1:17 (GeoSection) to 8:9 (AreaSection).
New York moves from 1:25 to 13:13.
Florida moves from 1:27 to 14:14.

We make three contributions:

\begin{enumerate}
\item \textbf{The AreaSection algorithm} (\S\ref{sec:algorithm}): dual
  population+area METIS bisection with isoperimetric ratio normalisation and
  a Lorenz-curve feasibility pre-filter.

\item \textbf{50-state empirical sweep} (\S\ref{sec:evaluation}): natural
  ratio results for all 50 states, showing which states are forced to near-equal
  splits and which retain asymmetric splits due to geographic concentration.

\item \textbf{Partisan comparison} (\S\ref{sec:evaluation}): head-to-head
  GeoSection vs. AreaSection on Wisconsin and North Carolina, showing identical
  seat allocations but different competitiveness profiles.
\end{enumerate}

\paragraph{Scope and relationship to B.8.}
AreaSection is a companion paper to GeoSection (B.8).
Both algorithms use ratio scanning with isoperimetric normalisation.
The difference is the vertex weight model: GeoSection uses ncon=1
(population only); AreaSection uses ncon=2 (population + land area).
All other algorithmic infrastructure --- recursive bisection tree,
orientation-aware subgraph splitting, proportionality analysis --- is shared.

AreaSection does not incorporate partisan data.
It is a purely geographic constraint: land area is a federal public record
(TIGER ALAND field), computed mechanically, with no interpretive degrees of
freedom.
\emph{We minimised edge-cut subject to geographic balance and verified the
result is not implicitly partisan.} This is a stronger and more legally durable
claim than minimising a partisan metric directly.
